\section{The \texttt{Flex\_monitor\_2D} McStas Component}
Flexible monitor.

\subsection*{Identification}
\begin{itemize}
  \item \textbf{Author:} Erik B Knudsen \& Peter Willendrup
  \item \textbf{Origin:} DTU Physics
  \item \textbf{Date:} Oct '20
\end{itemize}

\subsection*{Description}
Non-propagating 2D monitor that measures intensity (or something else) as a function of two selectable variables or parameters.

Example: Flex\_monitor\_2D(nU1=20, nU2=20, filename="Output", ustring1="x", ustring2="y", Umin1=-.1, Umax1=.1, Umin2=-.1, Umax2=.1)

\subsection*{Input parameters}
Parameters in \textbf{boldface} are required; the others are optional.

\begin{longtable}{p{0.22\textwidth}p{0.12\textwidth}p{0.46\textwidth}p{0.14\textwidth}}
\toprule
\textbf{Name} & \textbf{Unit} & \textbf{Description} & \textbf{Default} \\
\midrule
\endhead
nU1 & 1 & Number of U1 channels & 20 \\
nU2 & 1 & Number of U1 channels & 20 \\
filename & string & Name of file in which to store the detector image & 0 \\
\textbf{Umin1} &  & Minimum U1 to detect &  \\
\textbf{Umax1} &  & Maximum U1 to detect &  \\
uid1 & 1 & Integer index of uservar to be monitored. Overrides ustring1. & -1 \\
ustring1 & string & Name of variable U1 (user or particle state parameter as a string) to be monitored. & "" \\
\textbf{Umin2} &  & Minimum U1 to detect &  \\
\textbf{Umax2} &  & Maximum U1 to detect &  \\
uid2 & 1 & Integer index of uservar to be monitored. Overrides ustring2. & -1 \\
ustring2 & string & Name of variable U2 (user or particle state parameter as a string) to be monitored. & "" \\
signal & string & Name of variable to be used as an additive signal to be monitored. Default is intensity. & "p" \\
nowritefile & 1 & Flag to indicate if monitor should not save any data. & 0 \\
\bottomrule
\end{longtable}

\subsection*{Links}
\begin{itemize}
  \item Component source code found in file \texttt{Flex\_monitor\_2D.comp}.
\end{itemize}
\IfFileExists{monitors/Flex_monitor_2D_static.tex}{\input{monitors/Flex_monitor_2D_static.tex}}{}